% DSECS — Formal Theorem Appendix
% NeurIPS 2026 Submission
% \documentclass{article}
% \usepackage{amsmath,amssymb,amsthm}

\section{DSECS Formal Theorems}

\subsection{System Definition}

Let the system $S$ be defined as:

\[
S = (X, f_{\text{CSCO}}, f_{\text{GSAI}}, L)
\]

where:

\begin{itemize}
    \item $X$: input state space
    \item $f_{\text{CSCO}}: X \to \{0, 1\}$: constraint function
    \item $f_{\text{GSAI}}: X \to \mathbb{R}$: bounded utility scoring
    \item $L$: append-only hash-chain ledger
\end{itemize}

\newtheorem{theorem}{Theorem}
\newtheorem{corollary}{Corollary}
\newtheorem{lemma}{Lemma}

\begin{theorem}[Deterministic Execution]
\label{thm:determinism}
For all $x \in X$, $\text{execution}(x)$ is deterministic:
\[
\forall x: \text{run}(S, x) = y \text{ is unique}
\]
\end{theorem}

\begin{proof}
The system constructs a pure transition graph. Every transition applies $f_{\text{CSCO}}$ then $f_{\text{GSAI}}$ with deterministic hash-chain commitment through $L$. No random variables exist — the transition function accepts a required seed parameter. No hidden state mutation exists — \texttt{State} is an immutable dataclass with SHA-256 hash computed at initialization. The \texttt{clone()} method explicitly produces new instances with linked \texttt{parent\_hash}. By construction, identical input produces identical state trajectory across arbitrary repetitions. \end{proof}

\begin{theorem}[Constraint Safety]
\label{thm:safety}
For all $x \in X$:
\[
f_{\text{CSCO}}(x) = 0 \implies \text{execution}(x) \text{ is blocked}
\]
\end{theorem}

\begin{proof}
The runtime function \texttt{step()} is the single entry point for state evolution. When $\text{CSCO}(x) = \text{False}$, \texttt{step()} returns the current state unmodified. No alternative execution path exists — \texttt{step()} is the only function that advances system state. Therefore unsafe states are unreachable by construction. \end{proof}

\begin{theorem}[Bounded Utility]
\label{thm:bounded}
For all $x \in X$:
\[
f_{\text{GSAI}}(x) \in [0, 1]
\]
\end{theorem}

\begin{proof}
$f_{\text{GSAI}}$ is defined as:
\[
f_{\text{GSAI}}(x) = \text{stability} - 0.5 \times \text{failure\_rate} + 0.3 \times \text{throughput}
\]
where all input fields are bounded: $\text{stability} \in [0,1]$, $\text{failure\_rate} \in [0,1]$, $\text{throughput} \in [0, +\infty)$. All field values are clamped via \texttt{\_clamp} after every transition. The linear combination of bounded fields with finite coefficients produces bounded output. \end{proof}

\begin{theorem}[Ledger Consistency]
\label{thm:replay}
For all execution traces $T$:
\[
\text{Replay}(T) = T
\]
\end{theorem}

\begin{proof}
The ledger $L$ stores states as a hash-linked chain. Each append verifies $\text{state.parent\_hash} = \text{previous.hash}$. The replay function reconstructs each state using the stored transition seed and verifies hash equality at every step. Any mutation of ledger history is detected at the divergence point. \end{proof}

\begin{theorem}[Byzantine Robustness]
\label{thm:byzantine}
For $n = 3$ nodes with $f = 1$ adversarial node:
\[
\text{consensus output is invariant under single-node corruption}
\]
\end{theorem}

\begin{proof}
With $n=3$ and $f=1$, the honest majority (2 nodes) produces identical states under identical input and seed (by Theorem~\ref{thm:determinism}). The majority function selects the state with the most common hash. The single Byzantine node's output is outvoted 2:1. Tie-breaking uses lexicographically smallest hash among tied leaders. \end{proof}

\begin{corollary}[Full System Property]
\label{cor:full}
DSECS is a closed deterministic constrained decision system satisfying Theorems~\ref{thm:determinism} through~\ref{thm:byzantine} simultaneously.
\end{corollary}

\begin{lemma}[Hash Chain Injectivity]
\label{lem:hash}
The mapping from state content to hash is injective with overwhelming probability.
\end{lemma}

\begin{proof}
SHA-256 collision resistance ensures $P(\text{collision}) \leq 2^{-256}$. \end{proof}
